\section{Target Sum}
\label{sec:dp-target-sum}

\problemstatement{Given an array $\textit{nums}$ and an integer
$\textit{target}$, assign either a $+$ or $-$ sign to each element, and
count the number of distinct sign assignments that make the total
expression evaluate to exactly $\textit{target}$.}

\begin{patternbox}
\emph{``Assign $+$ or $-$ to each element to reach a target sum''} is,
word for word, Section~\ref{sec:dp-count-partitions-diff}'s partition
problem in disguise: every element assigned $+$ belongs to one ``subset''
$P$ (positive sum), every element assigned $-$ belongs to the other
subset $N$ (negative sum), and the expression's value is exactly
$P - N$. Counting sign assignments reaching $\textit{target}$ is
therefore exactly counting partitions with difference
$d = \textit{target}$ -- \textbf{no new derivation is needed at all.}
\end{patternbox}

\subsection*{Understanding the reduction}

Every element of $\textit{nums}$ goes into exactly one of two groups:
those assigned $+$ (summing to $P$) and those assigned $-$ (summing to
$N$, as a positive magnitude). The expression's total value is $P-N$,
and $P+N$ is always the array's total sum (every element contributes its
magnitude to exactly one of the two groups). This is precisely
Section~\ref{sec:dp-count-partitions-diff}'s setup, with $d =
\textit{target}$: count subsets summing to $(\textit{totalSum} +
\textit{target})/2$.

\subsection*{All Four Approaches}

Every stage is Section~\ref{sec:dp-count-partitions-diff}'s
implementation, called unchanged with $d = \textit{target}$.

\begin{lstlisting}
#include <bits/stdc++.h>
using namespace std;

int countSubsetsSpaceOptimized(vector<int>& arr, int k) {
    int n = arr.size();
    vector<int> prev(k + 1, 0);

    prev[0] = (arr[0] == 0) ? 2 : 1;
    if (arr[0] != 0 && arr[0] <= k) prev[arr[0]] = 1;

    for (int index = 1; index < n; index++) {
        vector<int> current(k + 1, 0);
        for (int target = 0; target <= k; target++) {
            int notPick = prev[target];
            int pick = (arr[index] <= target) ? prev[target - arr[index]] : 0;
            current[target] = notPick + pick;
        }
        prev = current;
    }
    return prev[k];
}

int findTargetSumWays(vector<int>& nums, int target) {
    int totalSum = accumulate(nums.begin(), nums.end(), 0);
    if ((totalSum + target) % 2 != 0 || totalSum + target < 0) return 0;

    int p = (totalSum + target) / 2;
    return countSubsetsSpaceOptimized(nums, p);
}

int main() {
    vector<int> nums = {1, 1, 1, 1, 1};
    cout << findTargetSumWays(nums, 3) << endl; // 5
    return 0;
}
\end{lstlisting}

\subsection*{Dry run}

\dryrun{Take $\textit{nums} = [1,1,1,1,1]$, $\textit{target}=3$.}

$\textit{totalSum}=5$. $P = (5+3)/2 = 4$. Counting subsets of five
$1$s summing to $4$: there are $\binom{5}{4}=5$ ways to choose which
four of the five $1$s are positive (the remaining one is negative),
giving $5$ valid sign assignments -- each evaluates to
$1+1+1+1-1=3=\textit{target}$, confirming the answer.

\begin{complexitybox}
\textbf{Time Complexity.} $O(n \cdot \textit{totalSum})$, inherited
from Section~\ref{sec:dp-count-subsets-sum-k}.
\textbf{Space Complexity.} $O(\textit{totalSum})$ for the rolling row.
\end{complexitybox}

\begin{takeawaybox}
This section is intentionally almost content-free beyond the reduction
itself: once Count Partitions with a Given Difference is solved, Target
Sum requires recognizing the disguise and nothing more. A surprising
number of seemingly distinct DP problems on interview sheets are exactly
this -- the same two or three core recurrences, wearing different cover
stories.
\end{takeawaybox}
